\documentclass{resume}

\begin{document}

\name{Harrison Chen}

\info{email}{hchen.robotics@gmail.com}
\info{phone}{+1(734)968-5295}
\info{location}{New York, NY}
\urlinfo{linkedin}{chen-harrison}{https://linkedin.com/in/chen-harrison}
\urlinfo{github}{chen-harrison}{https://github.com/chen-harrison}
\urlinfo{website}{chen-harrison.github.io}{https://chen-harrison.github.io}

\begin{personalstatement}
  Robotics software engineer with 5+ years experience developing autonomous navigation
  and perception systems for UAVs and automotive vehicles. Proficient in C++, Python, and ROS 2.
\end{personalstatement}

\section{Experience}

\entry{Senior Robotics Software Engineer}{OKSI}{Torrance, CA (Remote)}{Mar 2025 – Present}
\begin{details}
  \item Iterated on VSLAM GPS-denied navigation software product to increase robustness; improved accuracy
  of position estimate when recovering from temporary outages, added option to bypass FCU EKF when
  receiving input vehicle position, and optimized flight data recording to reduce CPU overhead while
  ensuring persistence on power loss
  \item Supported integration of navigation software into customer systems by debugging driver software,
  tailoring adapter nodes to handle incoming vehicle telemetry, and fixing edge cases captured by
  recorded flight data
  \item Implemented TensorRT inference and SAHI for YOLOv8 object detection, reducing model latency and
  increasing detection distance for high-resolution imagery
  \item Designed and maintained CI pipelines that perform code formatting checks, ROS bag playback testing,
  Docker image builds, and automated deployments
\end{details}

\entry{Autonomy Engineer}{PDW}{New Rochelle, NY}{Mar 2022 – Jan 2025}
\begin{details}
  \item Integrated GCOPTER trajectory generation algorithm into quadcopter’s autonomy stack, augmenting
  with BFS for selecting safe start and goal positions
  \item Evaluated depth estimation models HITNET and ESS DNN as a potential alternative to LiDAR;
  calibrated stereo Arducam sensors and benchmarked TensorRT and ONNX inference on Nvidia Jetson
  Xavier NX
  \item Wrote a ROS 1 package that supported OctoMap and Voxblox for creating a unified 3D occupancy grid
  map from forward- and downward-facing RealSense cameras
\end{details}

\entry{Robotics Engineer}{Jugaad Labs}{Philadelphia, PA}{Mar 2021 – Mar 2022}
\begin{details}
  \item Developed an automotive situational awareness system for semi-trucks, using CenterTrack to detect
  and track nearby vehicles in an image, then projecting the tracks into 3D space and applying a
  Kalman filter to the trajectories
  \item Built an application with Nvidia Isaac SDK to perform object detection in Isaac Sim warehouse
  environment, demonstrating a sensing architecture for autonomous logistics in simulation
\end{details}

% \entry{Applied Product Development Intern}{FANUC America}{Rochester Hills, MI}{Jun 2020 – Aug 2020}
% \begin{details}
%   \item Strengthened functionality for ArcTool recovery mechanism using proprietary programming language
%   Karel, enabling welding robots to recalibrate in any reachable end effector pose
% \end{details}

% \entry{Student / Team Member}{Robotic Systems Laboratory Course (ROB 550)}{Ann Arbor, MI}{Jan 2020 – May 2020}
% \begin{details}
%   \item Implemented a simulated SLAM robot in C++ with 2D occupancy grid mapping, odometry motion model,
%   beam measurement model, Monte Carlo localization, and A* path planning
%   \item Collaborated with teammates on the development of an inverted pendulum robot using C and RCL,
%   including PID control for balancing, manual steering via joystick, and autonomous movement along
%   series of waypoints
% \end{details}

\section{Education}

\entry{Master of Science in Robotics}{University of Michigan}{Ann Arbor, MI}{Dec 2020}
\begin{details}
  \item \textbf{GPA: 3.96/4.00}
  \item Relevant coursework: Mobile Robotics, Deep Learning for Computer Vision, Robot Modeling and Control
\end{details}

\entry{Bachelor of Science in Mechanical Engineering}{Northwestern University}{Evanston,
  IL}{June 2019}
\begin{details}
  \item \textbf{GPA: 3.80/4.00}
  \item Relevant coursework: Intro to Mechatronics, Machine Dynamics, Advanced Solid Modeling
  \item Activities: Education Chair @ Refresh Dance Crew, Social Chair @ Chinese Students Association, Tau
  Beta Pi
\end{details}

\section{Skills \& Interests}
\textbf{Software:} C++, Python, Bash, Git, Docker, MATLAB, CI/CD \\
\textbf{Robotics:} ROS 1/ROS 2, Eigen, OpenCV, PyTorch, Nvidia Jetson, sensor calibration \\
\textbf{Interests:} soccer, running, cooking, environmental conservation \\~\\

\textsuperscript{\textasteriskcentered} US Citizen, no work authorization required

\end{document}
